\appendixchapter{Foundations: Natural Numbers, Induction, and Choice}
\label{app:foundations}
\chaptermark{Natural Numbers, Induction, and Choice}
\section*{Introduction}
Chapter~\ref{ch:foundations} uses natural numbers, induction, and finite
recursive constructions directly. This optional appendix separates their
logical roles and gives worked guidance in induction before turning to
the genuinely different issue of choice. Throughout,
\[
\boxed{\N=\{1,2,3,\ldots\},\qquad 0\notin\N,\qquad \N_0=\N\cup\{0\}.}
\]
No particular set-theoretic realization of the natural numbers is needed.

Sections A.1--A.2 of Appendix~\ref{app:foundations} are a companion to
Chapter~\ref{ch:foundations}. They review induction, recursive definitions
and proofs, factorials, binomial coefficients, Pascal's identity, the finite
binomial theorem, and the equivalence of induction and well-ordering of $\N$.
Readers unfamiliar with these tools may consult those sections alongside
Chapter~2. Elementary finite algebraic identities such as the binomial
theorem may be used in the main text without proof; Appendix~A gives proofs
by induction for readers who want the foundational details.

Sections A.3--A.5 are fully optional foundations. Neither the equivalence
of Choice, Zorn's Lemma, and the Well-Ordering Theorem nor its transfinite
recursion and Hartogs arguments is needed for Chapters~3--9. For the main
text, knowing the statement and a typical use of Zorn's Lemma is more than
enough: if every chain in a nonempty partially ordered set has an upper
bound in that set, then the poset has a maximal element. A standard use is
the existence of a basis of an arbitrary vector space. The detailed
foundational proofs may be read purely for interest.


\section{The Natural Numbers and the Peano Axioms}
The \textbf{Peano axioms} express the elementary structure of \(\N\). At
this stage \(n+1\) denotes the successor of \(n\); the notation does not yet
presuppose that the usual binary addition operation has already been defined.
They assert:
\begin{enumerate}
\item \(1\in\N\);
\item if \(n\in\N\), then \(n+1\in\N\);
\item \(n+1\ne1\) for every \(n\in\N\);
\item if \(m+1=n+1\), then \(m=n\);
\item if \(A\subseteq\N\) satisfies
\[
1\in A,\qquad n\in A\Longrightarrow n+1\in A\qquad(n\in\N),
\]
then \(A=\N\).
\end{enumerate}
The final assertion is induction. We need neither model-theoretic
qualifications nor a separate successor function for these analysis notes.

\section{Addition, Induction, and Well-Ordering}
\subsection*{Addition by recursion}
The recursion principle permits a value at the next stage to be specified
from values already constructed.
\begin{theorem}[Recursion principle; stated without proof]
\label{thm:natural-recursion}
Let $X$ be a set, let $x_0\in X$, and let
\[
\Phi:\N_0\times X\longrightarrow X.
\]
There is a unique function $f:\N_0\to X$ such that
\[
f(0)=x_0,
\qquad
f(n+1)=\Phi(n,f(n))
\quad(n\in\N_0).
\]
\end{theorem}

Before using induction from zero, record the required form explicitly.
\begin{proposition}[Induction on \(\N_0\)]
\label{prop:n-zero-induction}
If $P(0)$ holds and
\[
P(n)\Longrightarrow P(n+1)
\qquad(n\in\N_0),
\]
then $P(n)$ holds for every $n\in\N_0$.
\end{proposition}
\begin{proof}
Apply ordinary induction on \(\N\) to the assertion $Q(k)=P(k-1)$.
Its base case is $P(0)$, and its induction step is exactly the displayed
implication with $n=k-1$. Thus $P(n)$ holds for every $n\geq0$.
\end{proof}

For each fixed \(m\in\N_0\), apply recursion with $X=\N_0$, initial
value $m$, and update rule $\Phi(n,x)=x+1$, where $+1$ still means
successor. This produces the unique map $n\mapsto m+n$ satisfying
\[
m+0=m,
\qquad
m+(n+1)=(m+n)+1.
\]
The final \(+1\) on the right initially means taking the successor. Once
this construction is made for every $m$, it defines a binary operation
\[
+:\N_0\times\N_0\longrightarrow\N_0.
\]
The earlier successor notation is then compatible with ordinary addition.
Induction now proves properties of the operation. For example, \(0+n=n\):
it holds at \(n=0\), and
\[
0+(n+1)=(0+n)+1=n+1.
\]
Thus
\[
\boxed{\text{recursion defines;}\qquad\text{induction proves}.}
\]

\begin{proposition}[Elementary addition laws]
For all $m,n,r\in\N_0$,
\[
0+n=n,
\qquad
m+0=m,
\qquad
(m+1)+n=(m+n)+1,
\]
and
\[
(m+n)+r=m+(n+r),
\qquad
m+n=n+m.
\]
\end{proposition}
\begin{proof}
The first identity was just proved; the second is part of the recursive
definition. Fix $m$ and induct on $n$ for successor compatibility. At
$n=0$, both sides equal $m+1$. If the identity holds at $n$, then
\[
(m+1)+(n+1)=((m+1)+n)+1=((m+n)+1)+1=(m+(n+1))+1.
\]

For associativity, fix $m,n$ and induct on $r$. The case $r=0$ is
immediate. If the identity holds at $r$, then
\[
\begin{aligned}
(m+n)+(r+1)
&=((m+n)+r)+1\\
&=(m+(n+r))+1\\
&=m+(n+(r+1)).
\end{aligned}
\]
Finally fix $m$ and induct on $n$ for commutativity. The case $n=0$
uses $0+m=m=m+0$. If $m+n=n+m$, then successor compatibility gives
\[
m+(n+1)=(m+n)+1=(n+m)+1=(n+1)+m.
\]
Each proof uses the previously established law and does not smuggle the
desired arithmetic into the recursion principle.
\end{proof}

Further arithmetic---multiplication, order, \(\Z\), and \(\Q\)---is
developed in Chapter~\ref{ch:foundations} rather than duplicated here.

\subsection*{How an induction proof works}
To prove \(P(n)\) for every \(n\in\N\), organize the argument as
\[
\boxed{\text{state }P(n)\to\text{base case}\to\text{induction hypothesis}
\to\text{prove }P(n+1)\to\text{conclude}.}
\]
The induction hypothesis is not an assumption that the theorem is already
true. It is the temporary assumption that \(P(n)\) holds at one arbitrary
stage \(n\), used to prove the next stage.

\begin{example}[A finite-sum identity]
For every \(n\in\N\),
\[
1+2+\cdots+n=\frac{n(n+1)}2.
\]
\textbf{Base case.} At \(n=1\), both sides equal \(1\).

\textbf{Induction hypothesis.} Fix arbitrary \(n\in\N\) and suppose
\[
1+2+\cdots+n=\frac{n(n+1)}2.
\]
\textbf{Induction step.} Add the next term and use the hypothesis:
\begin{align*}
1+2+\cdots+n+(n+1)
&=\frac{n(n+1)}2+(n+1)\\
&=\frac{(n+1)(n+2)}2.
\end{align*}
\textbf{Conclusion.} By induction, the identity holds for every
\(n\in\N\).
\end{example}

\begin{proposition}[Bernoulli's inequality]
If \(x\in\R\) and \(x\geq-1\), then, for every \(n\in\N_0\),
\[
(1+x)^n\geq1+nx.
\]
\end{proposition}
\begin{proof}
The assertion is equality at \(n=0\). If it holds at \(n\), then
\begin{align*}
(1+x)^{n+1}
&=(1+x)^n(1+x)\\
&\geq(1+nx)(1+x)\\
&=1+(n+1)x+nx^2\\
&\geq1+(n+1)x.
\end{align*}
Here \(1+x\geq0\) justifies multiplying the induction inequality.
\end{proof}

\subsection*{Common variants}
\paragraph{Shifted induction.}
If \(P(n_0)\) holds and
\[
P(n)\Longrightarrow P(n+1)\qquad(n\geq n_0),
\]
then \(P(n)\) holds for every \(n\geq n_0\). This is ordinary induction
with a shifted start and is useful when a statement is meaningful only after
an initial index.

\paragraph{Strong induction.}
Suppose \(P(1)\) holds and, for every \(n\in\N\),
\[
P(1),P(2),\ldots,P(n)\Longrightarrow P(n+1).
\]
Then \(P(n)\) holds for every \(n\in\N\). This is not a stronger axiom.
Define
\[
Q(n):\quad P(1),P(2),\ldots,P(n)\text{ all hold}.
\]
The initial case gives \(Q(1)\). If \(Q(n)\) holds, the strong-induction
hypothesis gives \(P(n+1)\), hence \(Q(n+1)\). Ordinary induction proves
every \(Q(n)\), and therefore every \(P(n)\). Thus
\[
\boxed{\text{when the next step needs more information, strengthen the induction statement}.}
\]

\begin{theorem}[Prime factorization: existence]
Every integer \(n\geq2\) can be written as a product of primes.
\end{theorem}
\begin{proof}
Use strong induction beginning at \(2\). The integer \(2\) is prime. Fix
\(n\geq2\) and suppose every integer \(k\) with \(2\leq k\leq n\) is a
product of primes. If \(n+1\) is prime, there is nothing to prove. If it is
composite, then
\[
n+1=ab,
\qquad
2\leq a,b<n+1.
\]
The induction hypothesis applies to both smaller factors. Writing each as a
product of primes and multiplying gives the required representation.
\end{proof}
Knowing the assertion only at \(n\) would not address the possibly much
smaller factors \(a\) and \(b\), which is why strong induction is natural.

\paragraph{Two-step induction.}
If
\[
P(n_0),\qquad P(n_0+1)
\]
hold, and
\[
P(n)\text{ and }P(n+1)\Longrightarrow P(n+2)
\qquad(n\geq n_0),
\]
then \(P(n)\) holds for every \(n\geq n_0\). Reduce this to ordinary
induction by defining
\[
Q(n):\quad P(n)\text{ and }P(n+1).
\]
The two initial cases give \(Q(n_0)\), and the two-step hypothesis turns
\(Q(n)\) into \(Q(n+1)\). This is convenient for recurrences depending on
the previous two stages. The same idea extends to any fixed finite number of
preceding stages.

\paragraph{Cauchy forward--backward induction.}
This different pattern proves a statement at powers of two and then steps
backward to every smaller index. Its classical model is the finite
arithmetic--geometric mean inequality. Let \(P_n\) denote
\[
\sqrt[n]{a_1\cdots a_n}\leq\frac{a_1+\cdots+a_n}{n},
\qquad a_i\geq0.
\]
We prove every \(P_n\), using the two-variable case
\cref{prop:two-variable-amgm}.

First, \(P_n\Rightarrow P_{2n}\). Apply \(P_n\) to each block of \(n\)
numbers and write their arithmetic means as \(A\) and \(B\). Then
\[
\sqrt[2n]{a_1\cdots a_{2n}}
\leq\sqrt{AB}\leq\frac{A+B}{2}
=\frac{a_1+\cdots+a_{2n}}{2n}.
\]
The first inequality follows by multiplying the two block inequalities and
taking the nonnegative square root; the second is two-variable AM--GM.

Second, \(P_n\Rightarrow P_{n-1}\) for \(n\geq2\). Given
\(a_1,\ldots,a_{n-1}\geq0\), put
\[
A=\frac{a_1+\cdots+a_{n-1}}{n-1}
\]
and append \(a_n=A\). The assertion \(P_n\) gives
\[
\sqrt[n]{a_1\cdots a_{n-1}A}\leq A.
\]
If \(A=0\), every \(a_i=0\) and \(P_{n-1}\) is immediate. If \(A>0\),
raise the inequality to the \(n\)th power and cancel \(A>0\):
\[
a_1\cdots a_{n-1}\leq A^{n-1}.
\]
Taking the nonnegative \((n-1)\)st root proves \(P_{n-1}\).

The forward steps give
\[
P_2\Longrightarrow P_4\Longrightarrow P_8\Longrightarrow\cdots.
\]
For a prescribed \(n\), choose a power \(2^k\geq n\), prove \(P_{2^k}\),
and apply the backward step until reaching \(P_n\). This completes the
equal-weight finite AM--GM proof. It is called Cauchy forward--backward
induction, not two-step induction: its two implications are
\(P_n\Rightarrow P_{2n}\) and \(P_n\Rightarrow P_{n-1}\).

\subsection*{Binomial coefficients and the binomial theorem}
Using the arithmetic developed in Chapter~\ref{ch:foundations}, define, for
\(n\in\N_0\) and \(0\leq k\leq n\),
\[
\binom nk=\frac{n!}{k!(n-k)!},
\qquad
0!=1.
\]
Adopt the extended convention
\[
\boxed{\binom nk=0\qquad\text{if }k<0\text{ or }k>n.}
\]
In particular, \(n<k\) implies
\[
\binom nk=0.
\]
 This convention simplifies
endpoint formulas and reindexing.

\begin{proposition}[Pascal's identity]
For \(n\in\N_0\) and every integer \(k\),
\[
\boxed{\binom n{k-1}+\binom nk=\binom{n+1}k.}
\]
\end{proposition}
\begin{proof}
For \(1\leq k\leq n\),
\begin{align*}
\binom n{k-1}+\binom nk
&=\frac{n!}{(k-1)!(n-k+1)!}+\frac{n!}{k!(n-k)!}\\
&=\frac{n!k+n!(n-k+1)}{k!(n-k+1)!}\\
&=\frac{(n+1)!}{k!(n+1-k)!}\\
&=\binom{n+1}k.
\end{align*}
At the endpoints and outside the range, the zero convention gives the same
identity without separate formulas.
\end{proof}

\begin{theorem}[Binomial theorem]\label{thm:appendix-binomial}
For \(x,y\in\R\) and \(n\in\N_0\),
\[
\boxed{(x+y)^n=\sum_{k=0}^{n}\binom nk x^{n-k}y^k.}
\]
\end{theorem}
\begin{proof}
The formula is immediate for \(n=0\). Assume it at \(n\). Multiply by
\(x+y\), split the two resulting sums, reindex, and use Pascal's identity:
\begin{align*}
(x+y)^{n+1}
&=\sum_{k=0}^{n}\binom nk x^{n+1-k}y^k
 +\sum_{k=0}^{n}\binom nk x^{n-k}y^{k+1}\\
&=\sum_{k=0}^{n+1}\binom nk x^{n+1-k}y^k
 +\sum_{k=0}^{n+1}\binom n{k-1}x^{n+1-k}y^k\\
&=\sum_{k=0}^{n+1}\left(\binom nk+\binom n{k-1}\right)
 x^{n+1-k}y^k\\
&=\sum_{k=0}^{n+1}\binom{n+1}k x^{n+1-k}y^k.
\end{align*}
Thus the mechanism is
\[
\text{multiply}\to\text{split}\to\text{reindex}\to\text{combine by Pascal}.
\]
\end{proof}
The proof uses only commutativity and distributivity, so the same identity
holds in any commutative ring. Ring theory is not needed for its use here.

\subsection*{Natural-number well-ordering}
\begin{proposition}[Well-ordering principle for \(\N\)]
\label{prop:natural-well-order}
Every nonempty subset of \(\N\) has a least element. In elementary
arithmetic this principle is equivalent to induction.
\end{proposition}
\begin{proof}
Assume induction and suppose \(E\subseteq\N\) has no least element. Prove
by induction that \(E\cap\{1,\ldots,n\}\) is empty. It is empty for \(n=1\);
if it is empty at \(n\), membership of \(n+1\) in \(E\) would make \(n+1\)
least. Thus \(E\) is empty, a contradiction.

Conversely, suppose a property holds at \(1\) and passes from \(n\) to
\(n+1\). If the failures formed a nonempty set, its least element \(m\)
could not be \(1\). Its predecessor \(m-1\) would satisfy the property,
forcing it at \(m\), a contradiction.
\end{proof}
This concerns the usual, already specified order on \(\N\), not the
Choice-equivalent Well-Ordering Theorem for arbitrary sets. This completes
the natural-number discussion; further arithmetic belongs to Chapter~2.

\section{The Axiom of Choice}
Finite or recursively specified constructions supply a rule for each next
step. Choice addresses a different problem: simultaneous selections from an
arbitrary family when no rule has been supplied.
\begin{definition}[Axiom of Choice]\label{def:axiom-choice}
For every indexed family \((X_i)_{i\in I}\) of nonempty sets, there is a
function
\[
c:I\longrightarrow\bigcup_{i\in I}X_i
\]
with \(c(i)\in X_i\) for every \(i\in I\).
\end{definition}
The sets need not be disjoint. A choice function for a finite family can be
constructed by induction, and an explicit rule needs no choice axiom.
Countable choice and dependent choice are weaker principles sufficient for
many familiar analytic selections.

\section{Choice, Zorn's Lemma, and Well-Ordering}
Recall from Chapter~\ref{ch:foundations} that a partially ordered set, or
poset, is a set $X$ equipped with a partial order $\leq$. Fix such a poset
$(X,\leq)$ throughout the following definitions. Two elements need not be
comparable; this is exactly why ``maximal'' and ``greatest'' differ.

\begin{definition}[Chains and upper bounds]
A subset $S\subseteq X$ is a \textbf{chain} if the restricted order is
total: for every $s,t\in S$, either $s\leq t$ or $t\leq s$. An element
$M\in X$ is an \textbf{upper bound} of $S$ if $s\leq M$ for every
$s\in S$. The upper bound need not belong to $S$.
\end{definition}
The empty set is a chain. If $X$ is nonempty, every element of $X$
vacuously upper-bounds it. This convention makes Zorn's statement uniform.

\begin{definition}[Maximal and greatest elements]
An element $x\in X$ is \textbf{maximal} if
\[
x\leq y\Longrightarrow y=x.
\]
An element $g\in X$ is \textbf{greatest} if
\[
x\leq g\qquad\text{for every }x\in X.
\]
\end{definition}
In the inclusion poset $X=\{\{1\},\{2\}\}$, both elements are maximal,
but neither is greatest because they are incomparable. If a poset is totally
ordered, however, a maximal element is greatest: comparability gives either
$x\leq g$ or $g\leq x$, and maximality of $g$ turns the second possibility
into equality. A maximal or greatest element need not exist.

\begin{theorem}[Zorn's Lemma]\label{thm:zorn}
If a nonempty partially ordered set has an upper bound for every chain, then
it has a maximal element.
\end{theorem}
\begin{theorem}[Well-Ordering Theorem]
\label{thm:well-ordering-theorem}
Every set admits a total order in which every nonempty subset has a least
element.
\end{theorem}

We need a small amount of transfinite language to prove equivalence rather
than merely announce it. A \textbf{well-order} is a total order in which
every nonempty subset has a least element. Two well-ordered sets have the
same \textbf{order type} if there is an order-preserving bijection between
them. An \textbf{ordinal} is the canonical order type of a well-ordered set;
its elements are precisely the earlier stages. At a successor stage one
uses the preceding value, while at a limit stage one uses the entire family
of earlier values.

\begin{theorem}[Transfinite recursion; stated without proof]
\label{thm:transfinite-recursion}
Let $\alpha$ be an ordinal. If a typed rule assigns an element of a set $Y$
to every function whose domain is an initial segment $\beta<\alpha$, then
there is a unique function $f:\alpha\to Y$ whose value at each $\beta$ is
the value prescribed from $f|_\beta$.
\end{theorem}
The construction is the same ``define from earlier stages'' mechanism as
ordinary recursion. Existence is proved by taking the union of all compatible
partial solutions; uniqueness follows by considering the first stage at which
two proposed solutions differ.

\begin{lemma}[Hartogs]
\label{lem:hartogs}
For every set $X$ there is an ordinal $h(X)$ for which no injection
$h(X)\to X$ exists.
\end{lemma}
\begin{proof}
Code a well-order on $A\subseteq X$ by the pair $(A,R)$, where
\[
R\subseteq A\times A\subseteq X\times X,\qquad
(A,R)\in\mathcal P(X)\times\mathcal P(X\times X).
\]
The pairs that code well-orders therefore form a set. Replace each by its ordinal
order type and let $\gamma$ be an ordinal strictly larger than every type
that occurs, for example the successor of their supremum. If $\gamma$
injected into $X$, its image, transported with the ordinal order, would be a
well-ordered subset of $X$ of type $\gamma$, contradicting the construction.
Take $h(X)=\gamma$.
\end{proof}

\begin{theorem}[Equivalence of choice principles]
\label{thm:choice-equivalence}
Over the usual axioms of set theory without choice,
\[
\boxed{
\mathrm{AC}\Longleftrightarrow\text{Well-Ordering Theorem}
\Longleftrightarrow\text{Zorn's Lemma}.}
\]
\end{theorem}
\begin{proof}
\textbf{AC $\Rightarrow$ Well-Ordering.}
Fix a set $X$. AC supplies a choice function
\[
c:\mathcal P(X)\setminus\{\varnothing\}\longrightarrow X,
\qquad c(A)\in A.
\]
To make stopping a typed recursion, adjoin a symbol $\dagger\notin X$
and use $Y=X\cup\{\dagger\}$ as target. Once $\dagger$ has appeared,
all later values are $\dagger$; when no unused point remains, the next
value is $\dagger$. Otherwise use the choice rule below. This defines a
rule on every initial history (ignore any $\dagger$ when forming its range).
Use transfinite recursion through the Hartogs ordinal $h(X)$: at stage $\alpha$, if
$X\setminus\{x_\beta:\beta<\alpha\}$ is nonempty, choose
\[
x_\alpha=c\bigl(X\setminus\{x_\beta:\beta<\alpha\}\bigr).
\]
If the process never exhausted $X$, the distinct values $x_\alpha$ would
define an injection $h(X)\to X$, contrary to \cref{lem:hartogs}. Thus it
has a first $\dagger$ stage $\lambda<h(X)$, and $\alpha\mapsto x_\alpha$ is a bijection
from $\lambda$ onto $X$. Transport the ordinal order to $X$; this well-orders
$X$.

\textbf{Well-Ordering $\Rightarrow$ Zorn.}
Let $P$ satisfy Zorn's hypotheses, and well-order its underlying set. Build
a chain by transfinite recursion through $h(P)$, using a stop symbol
$\dagger\notin P$ and target $P\cup\{\dagger\}$ as above. For any history
that is not a strictly increasing chain, prescribe $\dagger$ as well;
thus the rule is typed on every history. After selecting the chain
$C_\alpha=\{x_\beta:\beta<\alpha\}$, choose the least element of the fixed
well-order that is strictly above every member of $C_\alpha$, if such an
element exists; otherwise stop. The selected elements are distinct and form
a chain, so continuing without a stop would inject $h(P)$ into $P$,
contrary to Hartogs's lemma.
At the stopping stage let $C$ be the resulting chain. By hypothesis it has
an upper bound $u$. If $u\notin C$, then $u$ is strictly above every member
of $C$: upper-bound status gives $x\leq u$, distinctness gives $x\ne u$,
and antisymmetry rules out $u\leq x$. This contradicts the stopping rule;
hence $u\in C$. If $u<y$, then
$y$ is strictly above every member of $C$, another contradiction. Therefore
$u$ is maximal.

\textbf{Zorn $\Rightarrow$ AC.}
Let $(X_i)_{i\in I}$ be a family of nonempty sets. A partial choice function
is a function
\[
f:J\longrightarrow\bigcup_{i\in I}X_i,
\qquad J\subseteq I,
\qquad f(i)\in X_i.
\]
Order partial choice functions by extension. This poset is nonempty because
it contains the empty function. If $\mathcal C$ is a chain, then

\[
f=\bigcup_{g\in\mathcal C}g
\]
 is a function: two members that assign a value
at the same index are comparable by extension and hence agree. It is again a
partial choice function and upper-bounds the chain. Zorn gives a maximal
partial choice function
\[
f:J\to\bigcup_iX_i.
\]
 If $J\ne I$, choose one
$i_0\in I\setminus J$ and one $x_0\in X_{i_0}$. Then
$f\cup\{(i_0,x_0)\}$ is a larger partial choice function, a contradiction.
Thus $J=I$ and $f$ is a choice function. Selecting a single element from
the one already-given nonempty set $X_{i_0}$ in this contradiction does not
invoke the Axiom of Choice; the axiom concerns simultaneous choices from an
arbitrary family.
\end{proof}

The well-ordering principle for the already ordered set $\N$, proved in
\cref{prop:natural-well-order}, is equivalent to ordinary induction in this
elementary setting. It is not the Well-Ordering Theorem for arbitrary sets,
which the preceding theorem shows to be equivalent to AC.

\section{A Basis in Every Vector Space}
Recall that spans use finite linear combinations, even in infinite dimensions.
\begin{theorem}[Basis existence from Zorn's Lemma]\label{thm:zorn-basis}
Every vector space over a field has a basis.
\end{theorem}
\begin{proof}
Order the linearly independent subsets of \(V\) by inclusion. The poset is
nonempty because it contains \(\varnothing\). For a nonempty chain
\(\mathcal C\), let
\[
U=\bigcup_{E\in\mathcal C}E.
\]
 Every finite collection
of vectors of \(U\) lies in one member of the chain, so \(U\) is independent
and bounds the chain; the empty chain is bounded by \(\varnothing\). Zorn's
Lemma gives a maximal independent set \(B\). If
\(v\notin\operatorname{span}B\), then \(B\cup\{v\}\) is independent,
contradicting maximality. Hence \(B\) spans and is a basis.
\end{proof}

\Needspace{10\baselineskip}
\section*{Exercises}
\addcontentsline{toc}{section}{Exercises}
\begin{hexercise}{app-a-1}
For $a\in\R$, use recursion to define \(a^n\) for \(n\in\N_0\) and induction to prove
\(a^{m+n}=a^ma^n\). Identify which part is a definition and which is a proof.
\end{hexercise}
\begin{exercise}
Prove by ordinary induction that
\[
1+3+5+\cdots+(2n-1)=n^2.
\]
Mark the base case, induction hypothesis, and induction step.
\end{exercise}
\begin{hexercise}{app-a-3}
Use strong induction to prove that every integer \(n\geq2\) is either prime
or divisible by a prime. Explain why all smaller integers are useful.
\end{hexercise}
\begin{exercise}
Let \(u_0=0\), \(u_1=1\), and \(u_{n+2}=u_{n+1}+u_n\). Use two-step
induction to prove \(u_n<2^n\) for every \(n\geq1\).
\end{exercise}
\begin{exercise}
Use Pascal's identity and the zero-outside-range convention to derive
\[
\sum_{k=0}^n\binom nk=2^n.
\]
State whether you invoke the binomial theorem or reprove this special case.
\end{exercise}
\begin{exercise}
In the poset of proper subsets of \(\{1,2,3\}\), identify all maximal
elements and explain why none is a greatest element.
\end{exercise}
\begin{hexercise}{app-a-7}
Modify the basis-existence argument to prove that any linearly independent
subset extends to a basis. State exactly where Zorn's Lemma is used.
\end{hexercise}
